\section{长时 Rollout 与系统工程瓶颈}

前面的 GPU 与 RL slides 已经给出两个事实：模型计算需要规则的大 batch 才高效，而 Agent 轨迹却长度不一、等待外部环境且经常重试。本节把这一矛盾展开为调度、存储、通信和失败恢复问题，关注单位时间获得多少有效训练信号，而不是只看设备是否繁忙。

\subsection{为什么 Agent 训练很快变成系统问题}
课程指出，随着任务长度增长，训练瓶颈常从 ``算子吞吐'' 转向 ``异步调度、内存管理、通信开销''。Agentic rollouts 越长，等待 I/O 和环境反馈的时间占比越高，GPU 可能并未被充分利用。

\begin{knowledgebox}{长时 rollout 的资源形态}
\begin{itemize}
    \item 计算资源：策略前向与反向传播；
    \item 存储资源：轨迹缓存、日志、回放数据；
    \item 通信资源：多机同步、参数广播、环境 RPC；
    \item 调度资源：异步任务池、失败重试、超时控制。
\end{itemize}
\end{knowledgebox}

\subsection{内存层级与通信}
在长 rollout 中，轨迹、activation、模型状态和环境日志同时占用存储，跨卡同步又会阻塞采样或更新。本小节承接官方 slide 61--70，区分片上 SRAM、HBM 与跨设备链路，并说明为什么减少数据搬运、选择合适 sharding 往往比单纯增加理论 FLOPs 更有效。
课程后段对 GPU memory hierarchy 做了专门解释，强调 ``算得快'' 并不总是瓶颈，``拿数据慢'' 往往才是。共享内存、缓存、全局显存、跨卡通信都直接影响 Agent 训练效率。

\begin{importantbox}{工程结论}
当 rollout 进入长链路后，系统优化优先级往往应从 ``再堆算力'' 改为 ``降通信、降等待、降无效重试''。否则训练账单增加，但有效学习信号并未同步增长。
\end{importantbox}

\subsection{面向 Agent 的系统优化策略}
\begin{enumerate}
    \item 异步 rollout 与训练解耦，减少互相阻塞；
    \item 把高频短工具调用本地化缓存，降低 RPC 往返；
    \item 对轨迹做分级存储，热数据放快存，冷数据批量归档；
    \item 用失败类型路由替代统一重试，减少无效计算。
\end{enumerate}

\begin{warningbox}{系统优化中的错配风险}
若评估口径没同步更新，系统优化可能带来 ``看似更快但更差'' 的假象。例如吞吐提高了，但平均任务质量下降；或延迟降低了，但错误恢复率明显下滑。
\end{warningbox}

\subsection{本章小结}
长任务 Agent 的上限高度依赖系统工程能力。训练与系统是一体两面，不做系统优化会让算法增益被基础设施损耗掉。


